\subsection{连续函数集合的简单紧性}

在本节中，X 表示一个紧空间，\(\mathcal{C}_{s}(X)\) 表示 X 上以域 K（等于 R 或 C）为值的连续函数空间。空间 \(\mathcal{C}_{s}(X)\) 赋以 X 上的简单收敛拓扑。

命题 2. --- 设 D 是 X 的一个稠密子集，A 是空间 \(\mathcal{C}_{s}(X)\) 的一个子集。下列条件等价：

\begin{enumerate}
\def\labelenumi{(\roman{enumi})}
\item
  A 在 \(\mathcal{C}_s(\mathbf{X})\) 中相对紧。
\item
  从 \(\mathbf{A}\) 的元素的每个无穷序列，可以抽取一个在 \(\mathcal{C}_s(\mathbf{X})\) 中收敛的序列。
\item
  \(\mathbf{A}\) 的元素的每个无穷序列在 \(\mathcal{C}_s(\mathbf{X})\) 中有一个极限点。
\item
  设 \((f_n)_{n\in \mathbb{N}}\) 是属于 A 的函数序列，\((x_m)_{m\in \mathbb{N}}\) 是 D 的点的序列。若累次极限
\end{enumerate}

\[
\gamma = \lim _ {m \rightarrow \infty} \lim _ {n \rightarrow \infty} f _ {n} (x _ {m}), \quad \delta = \lim _ {n \rightarrow \infty} \lim _ {m \rightarrow \infty} f _ {n} (x _ {m})
\]

存在，则它们相等。此外，我们有 \(\sup_{f\in \mathbf{A}}|f(x)| < +\infty\)（对所有 \(x\in \mathbf{X}\)）。

\begin{enumerate}
\def\labelenumi{(\roman{enumi})}
\tightlist
\item
  \(\Rightarrow\) (ii)：设 \(\overline{\mathbf{A}}\) 是 \(\mathbf{A}\) 在 \(\mathcal{C}_s(\mathbf{X})\) 中的闭包。假设 \(\mathbf{A}\) 是紧的，并考虑函数序列 \(f_n \in \mathbf{A}\)（\(n \in \mathbf{N}\)）。设 \(\phi\) 是连续映射 \(x \mapsto (f_n(x))_{n \in \mathbf{N}}\)（从 \(\mathbf{X}\) 到可度量化空间 \(\mathbf{K}^{\mathbf{N}}\) 中的映射）。像 \(\mathbf{X}'\)（即 \(\mathbf{X}\) 在 \(\phi\) 下的像）是一个紧可度量化空间，因为 \(\mathbf{X}\) 是紧的。设 \(\mathbf{E}\) 是 \(\mathcal{C}_s(\mathbf{X})\) 的闭子空间，由满足下列条件的连续函数 \(f\)（定义在 \(\mathbf{X}\) 上）组成：关系 \(\phi(x) = \phi(y)\) 蕴含 \(f(x) = f(y)\)（\(x, y\) 是 \(\mathbf{X}\) 中任意一对点）。根据 GT, I, § 9, No.~4 的推论 2 与 GT, I, § 5, No.~2 的命题 3，映射 \(f' \mapsto f' \circ \phi\) 是一个同胚 \(\phi^*\)（从 \(\mathcal{C}_s(\mathbf{X}')\) 到 E 上）。因此，集 \(\mathbf{A}' = (\phi^*)^{-1}(\overline{\mathbf{A}})\) 在 \(\mathcal{C}_s(\mathbf{X}')\) 中是紧的，并且显然存在元素 \(f_n'\)（属于 \(\mathbf{A}'\)）使得 \(\phi^*(f_n') = f_n' \circ \phi\) 等于 \(f_n\)。
\end{enumerate}

由于 \(X'\) 是一个紧可度量化空间，存在一个可数稠密子集 \(D'\)（在 \(X'\) 中）(GT, IX, §2, No.~8, 命题 12 与 §2, No.~9, 命题 16)。设 \(T_{1}\)（相应地，\(T_{2}\)）是 \(A'\) 上由 \(D'\)（相应地，\(X'\)）上的简单收敛拓扑诱导的拓扑。那么 \(T_{1}\) 可度量化，\(T_{2}\) 紧且比 \(T_{1}\) 细，因而 \(T_{1}\) 与 \(T_{2}\) 一致；换句话说，\(A'\) 是 \(\mathcal{C}_{s}(X')\) 的一个紧可度量化子空间。从而，存在一个序列 \((f_{n_k}^{\prime})\)（从 \((f_n^{\prime})\) 中抽取），它收敛到某个元素 \(f^{\prime}\)（属于 \(\mathcal{C}_s(\mathbf{X}')\)）。因此，序列 \((f_{n_k})\) 收敛到 \(f = f' \circ \phi\)（在 \(\mathcal{C}_s(\mathbf{X})\) 中）。

\begin{enumerate}
\def\labelenumi{(\roman{enumi})}
\setcounter{enumi}{1}
\item
  \(\Rightarrow\) (iii)：这是显然的。
\item
  \(\Rightarrow\) (iv)：假设 A 的元素的每个无穷序列在 \(\mathcal{C}_s(\mathbf{X})\) 中有一个极限点。设 \(x \in \mathbf{X}\)。映射 \(\phi_x: f \mapsto f(x)\)（从 A 到 K 中的映射）是连续的。于是，\(\phi_x(\mathbf{A})\) 中的每个无穷序列有一个极限点；由于域 K（等于 R 或 C）可度量化，集 \(\phi_x(\mathbf{A})\) 在 K 中相对紧，从而有界。换句话说，我们有 \(\sup_{f \in \mathbf{A}} |f(x)| < \infty\)。
\end{enumerate}

设 \(f_n, x_m, \gamma\) 与 \(\delta\) 如 (iv) 所述。设 \(f\) 是序列 \((f_n)\) 在 \(\mathcal{C}_s(\mathbf{X})\) 中的一个极限点，设 \(x\) 是序列 \((x_m)\) 在紧空间 \(\mathbf{X}\) 中的一个极限点。对每个 \(m\)，映射 \(h \mapsto h(x_m)\)（从 \(\mathcal{C}_s(\mathbf{X})\) 到 \(\mathbf{K}\) 中）是连续的。根据假设，我们有 \(f(x_m) = \lim_{n \to \infty} f_n(x_m)\)，从而 \(\gamma = \lim_{m \to \infty} f(x_m)\)；由于 \(f: \mathbf{X} \to \mathbf{K}\) 连续且 \(x\) 是序列 \((x_m)\) 的一个极限点，我们得到 \(\gamma = f(x)\)。用类似的方式，我们证明等式 \(\delta = f(x)\)，从而 \(\gamma = \delta\)。

\begin{enumerate}
\def\labelenumi{(\roman{enumi})}
\setcounter{enumi}{3}
\tightlist
\item
  \(\Rightarrow\) (i)：假设数 \(f(x)\)（当 \(f\) 取遍 \(\mathbf{A}\) 时）所成的集在 \(\mathbf{K}\) 中有界（对所有 \(x\in \mathbf{X}\)）。这等价于假设闭包 \(\overline{\mathbf{A}}\)（即 \(\mathbf{A}\) 在乘积空间 \(\mathbf{K}^{\mathbf{X}}\) 中的闭包）是紧的 (GT, I, § 9, No.~5)。假设 \(\mathbf{A}\) 在 \(\mathcal{C}_s(\mathbf{X})\) 中不是相对紧的。这就意味着存在一个函数 \(u\in \overline{\mathbf{A}}\) 和一个点 \(a\in \mathbf{X}\)，使得 \(u\) 在 \(a\) 处不连续。于是存在实数 \(\varepsilon >0\)，使得在每个邻域 \(\mathrm{U}\)（\(a\) 的邻域）中，存在一点 \(x\) 满足 \(|u(x) - u(a)|\geqslant \varepsilon\)。
\end{enumerate}

我们将用归纳法构造点的序列 \((x_{n})_{n\in \mathbf{N}}\)（点取自 \(\mathbf{D}\)）与元素的序列 \((f_n)_{n\in \mathbf{N}}\)（元素取自 \(\mathbf{A}\)），使它们满足下列关系：

\[
\left| u (x _ {m}) - u (a) \right| \geqslant \varepsilon \quad \text { for } \quad m \geqslant 1;\tag{$(1)_{m}$}
\]

\[
\left| u (x _ {i}) - f _ {m} (x _ {i}) \right| \leqslant \frac {1}{m + 1} \text { for } 0 \leqslant i \leqslant m - 1;\tag{$(2)_{m}$}
\]

\[
\left| f _ {m} (x _ {i}) - f _ {m} (a) \right| \leqslant \frac {1}{i + 1} \quad \text { for } \quad 0 \leqslant m \leqslant i.\tag{$(3)_{m,i}$}
\]

我们取 \(x_0 = a\)，其中 \(f_0\) 是 A 中任意取的（集 A 非空，否则它在 \(\mathcal{C}_s(\mathbf{X})\) 中将相对紧）。设 \(n \geqslant 1\)，且 \(x_0, x_1, \ldots, x_{n-1}, f_0, f_1, \ldots, f_{n-1}\) 满足关系式 \((1)_m, (2)_m\)（对 \(1 \leqslant m < n\)）与 \((3)_{m,i}\)（对 \(0 \leqslant m \leqslant i < n\)）。由于 \(u\) 属于 \(\overline{\mathbf{A}}\)，存在 \(f_n \in \mathbf{A}\) 满足 \((2)_n\)。设 \(V_n\) 是所有 \(x \in X\) 中满足 \(|f_m(x) - f_m(a)| \leqslant \frac{1}{n+1}\)（对 \(0 \leqslant m \leqslant n\)）者所成的集。由于 \(f_n\) 连续，\(V_n\) 是 \(a\) 的一个邻域；选取一点 \(x_n\)（在 \(D \cap V_n\) 中）使得 \(|u(x_n) - u(a)| \geqslant \varepsilon\)，于是 \((1)_n\) 与 \((3)_{m,n}\) 得到满足。因此，构造可以继续进行。

由于 \(u(\mathbf{X})\) 是 \(\mathbf{K}\) 的一个紧子集，存在一个序列 \((y_k)\)（从 \((x_m)\) 中抽取），使得极限 \(\gamma = \lim_{k\to \infty}u(y_k)\) 存在。由 (2) \(_m\)，我们有 \(u(x_i) = \lim_{n\to \infty}f_n(x_i)\)（对所有 \(i\in \mathbf{N}\)），从而

\[
\gamma = \lim _ {k \rightarrow \infty} \lim _ {n \rightarrow \infty} f _ {n} (y _ {k}).\tag{4}
\]

另一方面，我们有 \(f_{n}(a) = \lim_{i\to \infty}f_{n}(x_{i})\)（由 (3) \(_{m,i}\)），从而 \(f_{n}(a) = \lim_{k\to \infty}f_{n}(y_{k})\)。由于 \(x_0 = a\)，由 (2) \(_{m}\) 推出 \(\lim_{n\to \infty}f_n(a) = u(a)\)。于是，

\[
u (a) = \lim _ {n \rightarrow \infty} \lim _ {k \rightarrow \infty} f _ {n} (y _ {k}).\tag{5}
\]

最后，由 \((1)_{m}\) 得到 \(|\gamma - u(a)| \geqslant \varepsilon\)，于是 \(\gamma \neq u(a)\)。这与断言 (iv) 矛盾；这样我们就证明了 (iv) 蕴含 (i)。

